\section*{第 \ref{ch:g12:randvar} 章 --- 随机变量与二项分布}

\begin{solution}{exo:g12:randvar:1}
$\E(X) = -0.3 + 0 + 0.8 + 0.5 = 1$。
$\E(X^2) = 0.3 \times 1 + 0 + 0.4\times4 + 0.1\times25 = 4.4$，由
柯尼希–惠更斯公式得 $\V(X) = 4.4 - 1 = 3.4$，且
$\sigma(X) = \sqrt{3.4} \approx 1.84$。
\end{solution}

\begin{solution}{exo:g12:randvar:2}
\emph{1.} 这 $10$ 个问题是独立的伯努利试验，其中
$p = \frac14$：$X \sim \mathcal B\left(10, \frac14\right)$，
$\E(X) = 2.5$，
$\sigma(X) = \sqrt{10 \times \frac14 \times \frac34} = \sqrt{1.875}
\approx 1.37$。

\emph{2.} $\P(X = 0) = \left(\frac34\right)^{10} \approx 0.056$；
\[
\P(X = 5) = \binom{10}{5}\left(\frac14\right)^5\left(\frac34\right)^5
\approx 0.058;
\quad
\P(X \geq 1) = 1 - \left(\tfrac34\right)^{10} \approx 0.944 .
\]
\end{solution}

\begin{solution}{exo:g12:randvar:3}
独立同分布的试验：$X \sim \mathcal B(400,\ 0.05)$，于是
期望理赔次数为 $\E(X) = 20$，且
\[
\sigma(X) = \sqrt{400 \times 0.05 \times 0.95} = \sqrt{19} \approx 4.4 .
\]
\end{solution}

\begin{solution}{exo:g12:randvar:4}
收到的付款 $P$ 满足
$\E(P) = \frac{5 + 6}{6} = \frac{11}{6}$，因此公平价格为
$c = \frac{11}{6} \approx 1.83$ 欧元。公平收益 $G = P - \frac{11}6$
取值为 $-\frac{11}{6}, \frac{19}{6}, \frac{25}{6}$，概率分别为
$\frac46, \frac16, \frac16$：
\[
\V(G) = \E(G^2)
= \frac{4 \times 121 + 361 + 625}{6 \times 36} = \frac{1470}{216}
= \frac{245}{36} \approx 6.8,
\qquad \sigma(G) \approx 2.6 .
\]
公平游戏的\emph{平均}收益为零，但其结果仍有波动：
公平并不等于无风险。
\end{solution}

\begin{solution}{exo:g12:randvar:5}
$X \sim \mathcal B(8,\ 0.7)$。
\[
\P(X = 6) = \binom86 (0.7)^6(0.3)^2 \approx 0.296 ;
\]
\[
\P(X \geq 6) = \P(6) + \P(7) + \P(8)
\approx 0.296 + 8(0.7)^7(0.3) + (0.7)^8
\approx 0.296 + 0.198 + 0.058 = 0.552 ;
\]
$\P(X \geq 1) = 1 - (0.3)^8 \approx 0.99993$。众数：
$(n+1)p = 6.3$，故最可能的取值为 $k^* = 6$
（\cref{exo:g12:randvar:7}）。
\end{solution}

\begin{solution}{exo:g12:randvar:6}
乘客到场人数为 $X \sim \mathcal B(104,\ 0.9)$；当
$X \geq 101$ 时航班超售：
\[
\P(X \geq 101) = \sum_{k=101}^{104}\binom{104}{k}(0.9)^k(0.1)^{104-k}
\approx 0.006 .
\]
多售 $4\%$ 的机票，仅使约
$0.6\%$ 的航班出现意外——这正是超售背后的经济学道理。
\end{solution}

\begin{solution}{exo:g12:randvar:7}
\[
\frac{\P(X = k+1)}{\P(X = k)}
= \frac{\binom{n}{k+1}}{\binom nk}\cdot\frac{p}{1-p}
= \frac{n - k}{k + 1}\cdot\frac{p}{1-p},
\]
利用 $\binom{n}{k+1} = \binom nk \frac{n-k}{k+1}$。该比值 $\geq 1$
当且仅当 $(n-k)p \geq (k+1)(1-p)$，亦即 $np - k p \geq k - kp + 1 - p$，当且仅当
$k \leq (n+1)p - 1$。故概率在
$k + 1 \leq (n+1)p$ 时严格递增，此后递减，最大值在
$k^* = \floor{(n+1)p}$ 处取得（与 $k^* - 1$ 同时取得，当
$(n+1)p$ 为整数时）。
\end{solution}

\begin{solution}{exo:g12:randvar:8}
\emph{1.} 首次出现正面在第 $k$ 次投掷，即前 $k-1$ 次为反面，随后出现正面：
概率为 $\left(\frac12\right)^{k-1}\cdot\frac12 = 2^{-k}$，其中
$1 \leq k \leq 10$。获胜的总概率为
$\sum_{k=1}^{10} 2^{-k} = 1 - 2^{-10} = \frac{1023}{1024}$（仅在
连续十次反面时游戏失败）。

\emph{2.} 期望收益：
\[
\sum_{k=1}^{10} 2^k \cdot 2^{-k} = \sum_{k=1}^{10} 1 = 10 \text{ 欧元}.
\]
若无上限，则和式 $\sum_{k\geq1} 1$ 发散，期望收益为
无穷，尽管游戏几乎总是只支付一小笔钱——这正是
著名的 \emph{圣彼得堡悖论}，说明仅凭期望
并不能衡量一个游戏的价值。
\end{solution}

\begin{solution}{pb:g12:randvar:1}
\textbf{1.} $X \sim \mathcal B\!\left(3, \frac16\right)$：
$\P(X = 0) = \frac{125}{216}$，
$\P(X = 1) = \frac{75}{216}$，
$\P(X = 2) = \frac{15}{216}$，
$\P(X = 3) = \frac{1}{216}$。

\textbf{2.} $\E(X) = 3 \times \frac16 = \frac12$；
$V(X) = 3 \times \frac16 \times \frac56 = \frac{5}{12}$。

\textbf{3.} 期望赔付 $2\,\E(X) = 1$ 欧元，相对于
$1$ 欧元的赌注：收益为 $0$ --- 恰好公平（罕见情形）。

\textbf{4.} 不放回抽取时，各次抽取不独立（第二张牌的
概率取决于第一张）：不是二项分布 ---
检查清单中的独立性一项不满足。有放回抽取时：
$n = 5$ 固定，$p = \frac14$ 恒定，各次抽取独立：
二项分布。

\textbf{5.} $\E(X) = 6$，$\sigma(X) = \sqrt{4.2} \approx
2.05$。$\E(S) = 5 \times 6 - 10 = 20$；
$\sigma(S) = 5\sigma(X) \approx 10.25$。

\textbf{6.} 固定 $n = 105$ 次试验（即票数），每次结果为
到场/未到场之一，且 $p = 0.9$ 相同，假定各次独立：
二项分布。坦白地说，家庭常常会
\emph{一起}误机，暴风雨会让整架飞机空座 --- 独立性
是模型的一次信念之跃，而相关的未到场使
尾部比二项分布所承诺的更厚。

\textbf{7.} $\E(X) = 94.5$；
$\sigma(X) = \sqrt{105 \times 0.9 \times 0.1} =
\sqrt{9.45} \approx 3.07$。

\textbf{8.} $\P(X \geq 101) = \sum_{k=101}^{105}
\binom{105}{k} 0.9^k\, 0.1^{105 - k}$。

\textbf{9.} $\approx 0.0167$：约每六十个航班中有一班
至少会有乘客被挤下飞机。

\textbf{10.} 每班航班中，$\E(\text{被挤下}) \approx 0.023$ 名乘客：
期望赔偿 $\approx 19$ 欧元 --- 相对于
$1\,000$ 欧元的额外收入。超售五张票
压倒性地有利可图；因此各家航空公司都这样做。

\textbf{11.} 由尾概率表可得：
$n = 105$：$1.7\,\%$；$n = 106$：$4.0\,\%$；$n = 107$：
$8.1\,\%$。符合要求的最大售票数为 $n = 106$ 张。

\textbf{12.} $z = \frac{100.5 - 94.5}{3.07} \approx 1.95$：
阈值位于均值之上两个标准差处，
而钟形曲线的经验法则（“上侧 $2\sigma$ 尾部
$\approx 2.5\,\%$”）与精确的 $1.7\,\%$ 很接近 ---
光滑曲线对二项分布的逼近，正是即将到来的
各章的主角。

\textbf{13.} $0.98^{20} + 20 \times 0.02 \times 0.98^{19}
\approx 0.94$：合格批次有 $94\,\%$ 的概率通过抽检。

\textbf{14.} $0.9^{20} + 20 \times 0.1 \times 0.9^{19}
\approx 0.39$：不合格批次仍有 $39\,\%$ 的
概率混过抽检。

\textbf{15.} 生产者风险：合格批次被拒收
（$\approx 6\,\%$）；消费者风险：不合格批次被接收
（$\approx 39\,\%$）。更大的样本（配以按比例设定的
接收阈值）可以同时降低两者 --- 代价是更多的
检验：质量也要计入预算。

\textbf{16.} $V\!\left(\frac Xn\right) = \frac{V(X)}{n^2} =
\frac{np(1-p)}{n^2} = \frac{p(1-p)}{n}$：标准差为
$\sqrt{\frac{p(1-p)}{n}}$。样本量扩大四倍，
噪声减半：十年级波动区间中的 $\frac{1}{\sqrt n}$ 规律，
现在得到了推导。

\textbf{17.} 上限把赔付封顶为 $2^{10} = 1024$
欧元，因此十轮中的每一轮恰好贡献 $1$ 欧元
的期望，即 $\E = 10$。无上限游戏的无穷
期望完全来自极端罕见的
天文数字般巨大的赔付；封顶则承认，没有哪家银行
会支付 $2^{50}$ 欧元。封顶游戏的一个公平票价：
$10$ 欧元 --- 悖论顿时不复存在。

\textbf{18.} 每张票的期望奖金：
$\frac{100 + 2 \times 50}{200} = 1$ 欧元，相对于 $2$ 欧元的
票价：利润率为 $50\,\%$ --- 是轮盘赌
$2.7\,\%$ 的十八倍。彩票之所以能够维持，在于其参与者
明知是在捐赠：“亏损”正是关键所在。

\textbf{19.} 期望理赔额：$1\,000 \times 0.01 \times
10\,000 = 100\,000$；保费 $130\,000$：期望利润
$30\,000$ 欧元。总理赔额的标准差：
$10\,000 \times \sqrt{1\,000 \times 0.01 \times 0.99}
\approx 31\,500$ 欧元 --- 一个普通的糟糕年份就会吞掉
整个期望利润。因此需要资本准备金、再保险，
以及远大于一千份保单的组合：保险公司
依靠大数定律营运，并针对其
收敛缓慢预留准备金。

\textbf{20.} 先查清单 --- 固定 $n$，$p$ 相同，
独立性\emph{被宣称成立} --- 否则根本不是二项分布。然后
依据 $\E \pm \sigma$ 定位：均值指明位置，偏差发出警示。
再为尾部定价：超售挤下乘客、不合格批次、毁灭性的年份，
都落在 $2\sigma$ 之外。须注意：独立性是建模者的
承诺，而非世界本身的性质；而期望恰恰忽略了
足以毁掉你的东西。规则手册缺失的内容 --- 频率
如何趋于稳定，以及波动呈什么形态 ---
正是接下来两章的主题。
\end{solution}
